\documentclass{article}
\usepackage{amsmath,amssymb}
\usepackage{xurl}
\usepackage[hidelinks]{hyperref}
% Shared commands for notes in docs/paper-gaps/.

\DeclareUnicodeCharacter{2080}{\ifmmode{}_{0}\else\textsubscript{0}\fi}
\DeclareUnicodeCharacter{2081}{\ifmmode{}_{1}\else\textsubscript{1}\fi}
\DeclareUnicodeCharacter{2082}{\ifmmode{}_{2}\else\textsubscript{2}\fi}
\DeclareUnicodeCharacter{2099}{\ifmmode{}_{n}\else\textsubscript{n}\fi}

\newcommand{\Lean}{\textsc{Lean}}
\ExplSyntaxOn
\NewDocumentCommand{\leanid}{m}{
  \ifmmode
    \textsf{\detokenize{#1}}
  \else
    \group_begin:
    \urlstyle{sf}
    \tl_set:Nn \l_tmpa_tl {#1}
    \tl_replace_all:Nnn \l_tmpa_tl {\_} {_}
    \exp_args:NV \path \l_tmpa_tl
    \group_end:
  \fi
}
\ExplSyntaxOff
\providecommand{\tr}{\operatorname{tr}}
\newcommand{\tnleanIssueBase}{https://github.com/LionSR/TNLean/issues/}
\newcommand{\tnleanPrBase}{https://github.com/LionSR/TNLean/pull/}
\newcommand{\tnleanDocBase}{https://sirui-lu.com/TNLean/docs/}
\newcommand{\ghissue}[1]{\href{\tnleanIssueBase#1}{GitHub issue \##1}}
\newcommand{\ghpr}[1]{\href{\tnleanPrBase#1}{PR \##1}}
\newcommand{\doclink}[1]{\href{\tnleanDocBase#1.html}{\path{#1}}}

% Dirac notation and the identity operator, as in blueprint/src/macros/common.tex.
% \providecommand keeps note-local definitions authoritative.
\usepackage{dsfont}
\providecommand{\ket}[1]{|#1\rangle}
\providecommand{\bra}[1]{\langle #1|}
\providecommand{\braket}[2]{\langle #1 | #2 \rangle}
\providecommand{\ketbra}[2]{|#1\rangle\!\langle #2|}
\providecommand{\Id}{\mathds{1}}
\providecommand{\one}{\mathds{1}}
\providecommand{\C}{\mathbb{C}}
\providecommand{\MN}[1]{M_{#1}(\C)}
\providecommand{\MD}{M_D(\C)}

% Verdict marker: \gapnote{<kind>}{<status>}. Kind is the policy
% classification (clarification, local-correction, scope-restriction,
% unfaithful, false-source, open-gap); status is open, wip (elimination
% underway), resolved, or historical. Machine-read by the site index;
% prints a verdict line.
\newcommand{\gapnote}[2]{\par\noindent\textbf{Verdict:} #1 (#2).\par}

\title{Finite Adaptive Local Channels and the Number of Measurement Rounds}
\author{}
\date{2026-10-04}
\begin{document}
\maketitle
\gapnote{scope-restriction}{open}

\section{Source operations}
Piroli, Styliaris and Cirac~\cite{Piroli2021LOCC}, paragraph ``State
transformations with QC and LOCC'', define one circuit followed by sequential
local measurements and local unitaries conditioned on earlier outcomes.
The same paragraph mentions the possible extension to several rounds of
local operations and classical communication. The later paragraph ``Phases
of matter'' allows a number of composed channels that is fixed independently
of the system size.

\section{Finite adaptive channels}
An onsite instrument has local Kraus operators $K_{i,j}$.
At each site they satisfy $\sum_jK_{i,j}^{\dagger}K_{i,j}=I$.
For an outcome tuple $J$, let $K_J=\bigotimes_iK_{i,J(i)}$.
If $\Psi_J$ is the channel selected by that tuple, averaging the outcomes gives
\begin{align}
 \rho&\longmapsto\sum_J\Psi_J(K_J\rho K_J^{\dagger}).
\end{align}
These are unnormalized operations. Their traces already contain the outcome
probabilities, so no division by probabilities is needed, even on branches
whose probability is zero. Complete positivity follows from the Kraus
operators of the continuations, composed with $K_J$. Trace preservation
follows from $\sum_JK_J^{\dagger}K_J=I$ and trace preservation of each
$\Psi_J$.\footnote{Formal declarations:
\leanid{QuantumCircuit.OnsiteChannel.feedforwardMap} and
\leanid{QuantumCircuit.OnsiteChannel.feedforwardMap_isKrausCPTP} in
\doclink{TNLean/Circuit/Channel/Feedforward}.}

A finite adaptive protocol is built from onsite channels, layers of
nearest-neighbor channels, and onsite instruments followed by
outcome-dependent continuations. Its depth bounds the number of
nearest-neighbor layers on every branch; onsite operations and classical
communication cost zero depth. Repeated instrument steps allow later
measurements and gates to depend on the complete history. Such protocols
are channels on arbitrary input operators. Composing protocols of depth at
most $T_1$ and $T_2$ has depth at most $T_1+T_2$.\footnote{Formal declarations:
\leanid{QuantumCircuit.IsAdaptiveChannelProtocol},
\leanid{QuantumCircuit.IsAdaptiveChannelProtocol.isKrausCPTP}, and
\leanid{QuantumCircuit.IsAdaptiveChannelProtocol.comp} in
\doclink{TNLean/Circuit/Measurement/AdaptiveConversion}.}

The existing fixed-basis measurement round has Kraus operators
$V_mP_mU$, where $U$ is its circuit, the orthogonal projections $P_m$ sum to
$I$, and $V_m$ is the product of conditional local corrections.
Its Kraus operators satisfy
\begin{align}
  \sum_m(V_mP_mU)^{\dagger}(V_mP_mU)&=I.
\end{align}
Thus its average map is a channel on arbitrary inputs.\footnote{Formal declarations:
\leanid{QuantumCircuit.MeasurementRound.map} and
\leanid{QuantumCircuit.MeasurementRound.map_isKrausCPTP} in
\doclink{TNLean/Circuit/Measurement/Channel}.}
The round is also a finite adaptive protocol of depth equal to its circuit
length: an onsite instrument measures the selected sites, then the outcome
selects the onsite correction.\footnote{Formal declaration:
\leanid{QuantumCircuit.MeasurementRound.adaptive} in
\doclink{TNLean/Circuit/Measurement/AdaptiveRound}.}

\section{Physical gate-cost boundary}
The adaptive predicate inherits the enlarged-site convention of the local
channel model. Onsite channels may change the local dimension without a
uniform bound, and later gates on a pair of enlarged sites still cost one
layer. The source distinguishes fixed-dimensional physical qudits from
onsite ancillas. In particular, the Supplement's proof of Proposition 2
(p.~8) bounds the entanglement increase using the physical dimension,
independently of the number of onsite ancillas.

Consequently, the formal quantum-depth bound is not yet a bound in the
source's resource model. A comparison needs an explicit simulation by
physical-qudit gates, with a depth overhead controlled uniformly in the
chain length. A possible restricted target bounds every intermediate
onsite dimension by one constant before quantifying over chain lengths,
then derives an explicit simulation overhead depending on that constant
and the physical dimension. No such simulation is established here.

\section{Asymptotic boundary}
The finite adaptive relation does not bound the number of measurement rounds
independently of system size. Its depth bound counts quantum layers only,
so a family may have arbitrarily many free instrument steps as the chain
length grows. It must therefore not be substituted for the source's
fixed-number-of-rounds asymptotic phase relation.
In addition to the physical gate-cost comparison, retain a separate bound on the number of composed
single-round channels and quantify that bound before the chain length.
The finite adaptive conversion results address the finite-round task; the
fixed-round asymptotic relation and MPS classification remain separate.
\footnote{Finite adaptive conversions are tracked by \ghissue{8382}; the
asymptotic relation and classification are tracked under \ghissue{8622}.}

\bibliographystyle{alpha}
\bibliography{references}
\end{document}
